% Auto-generated by scripts/generate_tables/top_states.py
% Source: research/analysis/moratorium_inventory.csv
\begin{figure}[tp]
\centering
\caption{Top States by Moratorium Count}
\label{tab:top-states}
\begin{tikzpicture}
\begin{axis}[
    xbar,
    bar width=10pt,
    width=0.78\textwidth,
    height=6.5cm,
    xlabel={Number of Instruments},
    xmin=0,
    xmax=160,
    ytick=data,
    symbolic y coords={
        Minnesota,
        Illinois,
        Indiana,
        Kentucky,
        Wisconsin,
        Washington,
        Tennessee,
        Iowa,
        Georgia,
        North Carolina,
        Michigan,
        Ohio
    },
    nodes near coords,
    nodes near coords style={font=\small\bfseries, anchor=west, color=black},
    every axis plot/.append style={fill=chart-bar, draw=chart-bar-dark},
    axis lines*=left,
    xmajorgrids=true,
    grid style={gray!20},
    tick label style={font=\small},
    label style={font=\small},
    y dir=normal,
]
\addplot[fill=chart-bar, draw=chart-bar-dark] coordinates {
    (26,Minnesota) (27,Illinois) (28,Indiana) (31,Kentucky) (33,Wisconsin) (35,Washington) (43,Tennessee) (54,Iowa) (69,Georgia) (70,North Carolina) (141,Michigan) (159,Ohio)
};
\end{axis}
\end{tikzpicture}
\begin{flushleft}
\small\textit{Note:} Counts reflect individual instruments in the cleaned
inventory (\Cref{app:inventory}, $n = 1051$). Each entry represents a unique
jurisdiction--instrument pair after removing federal entities, non-governmental
actors, and duplicate extensions of the same underlying moratorium.
\end{flushleft}
\end{figure}
